\section{Dimensional unity, relative symmetry, and recovery of structure}
\label{lib06:sintesis}

The holographic extreme and mean ratio articulates conservation of unity
under evolution with its arithmetic and dimensional realizations. The
composition developed here makes this principle precise through
conserved quantities, domains, and inverses. The enriched state
generated by APP--TRIT--TPK precedes the numerical and
geometric readouts; its operations preserve the orientation, carry,
incidence, and memory used by each realization.

\subsection{Arithmetic conservation, unitarity, and isometric decomposition}
\label{lib06:tres-conservaciones}

Arithmetic refinement preserves normalized unity:
\begin{equation}
 (x,y,c)\longmapsto(Bx-c,By+c),\qquad
 \frac{Bx-c}{B(x+y)}+\frac{By+c}{B(x+y)}=1.
 \label{lib06:unidad}
\end{equation}
The residue inverse recovers the canonical carry and the preceding
pair. Block composition preserves the contributions of
events with their scale, and successive divisions recover every
finite history. The limiting readout is controlled by decreasing
cylinders; boundary marks retain the distinction between
histories sharing a scalar evaluation.

The lift on labels transforms that bijection into a
unitary \(J\). It preserves the norm of amplitudes and transports
the fraction and carry readers through
\(J^*M_hJ=M_{h\circ\ell}\). The quantity \(x^2+y^2\) and the norm
of \(|x,y,c\rangle\) belong to different realizations. The
proof of unitarity establishes the precise link between arithmetic
reversibility and preservation of amplitudes.

The third conservation law corresponds to two isometric realizations:
\begin{equation}
 (a+1)Q_-^*Q_-+aQ_+^*Q_+
 =(2a+1)(a^2+a+1)I.
 \label{lib06:balance}
\end{equation}
Cancellation of the mixed terms allows bilateral
distribution and its norm-one reconstruction. The law holds for the
transports on the domain specified in
\cref{lib04:teo-balance}. Its specialization to the ternary cycle
provides additional norm and index properties; the general law
remains valid for isometries without that finite order.

The chain formed by the three maps is effective: the bijection
produces \(J\); a previously declared permutation produces the
second transport \(JR\); the two transports satisfy the
balance, and their channels recover the input through
\eqref{lib05:inversa-bilateral}. The example in
\cref{lib05:ejemplo} shows the complete procedure with
\((73,27,1)\), \((729,271)\), and \(1001\) admissible labels.
The numbers in that realization are not introduced as substitutes for
the structure generating the register \(K\).

\subsection{Complete memory and the logical contribution to the Landauer bound}
\label{landauer:puente-archivo}

The vanishing logical contribution to the Landauer bound, established in
\eqref{v:ant:eq:ii-kb-aut-reversible}, applies to the recoverable
transports constructed in this article. The bijection of
\cref{lib05:teo-biyeccion} preserves the carry through
\eqref{lib05:inversa-etiquetas}; its prolongation preserves the path
through \cref{lib05:teo-rutas}. The archives of
\cref{lib04:teo-archivo} preserve the amplitudes and correlations
needed to invert the transport.

\begin{corollary}[Recovery and spectral conservation of the archive]
\label{landauer:cor-archivo}
Let \(Z:H\to\mathcal K\) be one of the complete archives
\(Z_N,Z_\infty\) of \cref{lib04:teo-archivo}, or the unitary
lift \(J_n\) of \cref{lib05:teo-rutas}, on its respective domain.
For a finite-rank density operator \(\rho\) on \(H\),
the output \(\rho_Z=Z\rho Z^*\) satisfies
\begin{align}
 \operatorname{Tr}\rho_Z&=1,& Z^*\rho_ZZ&=\rho,\\
 s(\rho_Z)&=s(\rho),&
 s(\rho)&=-\operatorname{Tr}(\rho\log\rho),
 \label{landauer:espectro}
\end{align}
with \(0\log0=0\). If \(X\) is a random variable on the finite
paths of \cref{lib05:teo-rutas}, then
\begin{equation}
 \mathsf H(X\mid\Phi_n(X))=0.
 \label{landauer:recuperacion-ruta}
\end{equation}
\end{corollary}
\begin{proof}
Write the finite spectral decomposition
\(\rho=\sum_jp_j|v_j\rangle\langle v_j|\), with the \(v_j\)
orthonormal, \(p_j>0\), and \(\sum_jp_j=1\). The identity
\(Z^*Z=I\) implies that the \(Zv_j\) are orthonormal. Thus,
\(\rho_Z=\sum_jp_j|Zv_j\rangle\langle Zv_j|\) has the same
nonzero eigenvalues, proving the trace and entropy identities.
Multiplication by \(Z^*\) and \(Z\) recovers \(\rho\). Finally,
\(\Phi_n^{-1}\) recovers each path from its final pair; the
conditional distribution is concentrated on a single path and its
entropy is zero.
\end{proof}

Conservation pertains to the complete archive. A visible readout
\(q\), accompanied by memory \(M\), recovers the state under the
hypothesis of \eqref{v:ant:eq:ii-kb-aut-fibra}; the information not
expressed by \(q\) remains in
\(\mathsf H(M(X)\mid q(X))\). Discarding a complement or a
necessary terminal defines another operation, whose recoverability is
determined by the information actually retained. The preceding spectral
identity preserves the distinction between entropy of the complete
state, entropy of a partial reduction, and heat transferred.

In the thermal model of Section
\ref{v:ant:sec:ii-kb-aut-desarrollo}, with degenerate logical states,
an isothermal environment, and accounting for auxiliary registers, the
term associated with erasure is
\(k_B\Theta\,\mathsf H(X\mid F(X))\), for a deterministic logical
map \(F\). Its contribution vanishes for the preceding recoverable
transports; resetting a uniform TRIT satisfies
\(Q_{\rm reset}\geq k_B\Theta\log3\). Preparation, control,
finite operating speed, and coupling to the environment retain their
own physical balances. The total heat of an implementation requires
combining these contributions with the logical operation.

\subsection{Conservation of an observable and its relative symmetry}
\label{lib06:simetria}

Norm balance is the identity case of a more informative
operator law. When the transports are unitary, with relative operator
\(R\), the effective reader is
\begin{equation}
 \mathcal T_a(G)=\frac{a(a+1)G+R^*GR}{a(a+1)+1}.
 \label{lib06:lector}
\end{equation}
The observable remains fixed exactly when \([G,R]=0\).
In the canonical case,
\begin{equation}
 \mathcal T_8(G)-G=\frac{R^*GR-G}{73}.
 \label{lib06:defecto-lector}
\end{equation}
This identity determines both conservation and variation:
the defect measures the difference between two realizations of the same
reader, weighted by the split \(72+1\).

Complete memory also transports the observable through
\(A\mapsto ZAZ^*\) on \(\operatorname{im}Z\). Products,
adjoints, and pairings are preserved on that image. The diagonal
channel readout and the transported observable are different
operations; the carry example shows the difference between
\(74/73\) for the former and eigenvalue one for the latter.
The information remains recoverable even if a particular diagonal
readout changes value.

The variational charges of \cref{nuc11:noether} retain their
foundation in an explicit action and intertwining. For
\(A_{n+1}U_n=U_nA_n\), the transport action gives the charge
\(p_n^*A_nq_n\). Its conservation and the norm identity are
compatible results with different hypotheses. Likewise,
preservation of an alternating or Lorentzian form in
\eqref{lib04:balance-formas} retains its geometric meaning;
the noise and contraction estimates use the Hilbert
positivity declared in their proofs.

\subsection{Role of the register in transduction}
\label{lib06:registro}

The holographic arithmetic--geometric coupling constant has
a role defined by its compositions. The register \(K\)
is obtained beforehand by oriented and Hadamard extraction.
Its readout \(\kappa\) is reversible with the base, length, and origin
fixed. Its frame \(O_K\) identifies cyclic responses because
its Gram matrix has a strictly positive lower bound. Incidence
and geometric direction use their declared readers
on the same register.

The archive isometry preserves the Gram matrix of that frame and therefore
retains the recovery bounds at every depth. The
complete register admits norm-one inversion; the two partial
memories retain their factor \(9/4\) and their uniform kernel. The
additional charge distinguishes recovery of the absolute register from
that of its centred component. Integer decoding subsequently recovers
the incidence selection even at a threshold lying
exactly on a coordinate.

This capacity is more specific than preserving a total quantity:
it allows recovery of operations and relations acting on a
generated object. Its guarantees rest on the Gram matrix, the
distances in the integer image, and the proven inverses, not
on retaining another copy of the input by definition. At the same time,
the finite number of windows of \(K\) remains distinct from the
complete enriched history; the other registers and accompanying
marks retain their own transports.

The orbit and its isometric transports preserve the norm
\(\nu=\|K\|\). Normalizing its columns to norm one requires
retaining \(\nu\) to invert that operation. Polar normalization
uses the pair
\((F_K,G_K)\). The formula
\(O_K=F_KG_K^{1/2}\) preserves orientation and amplitudes.
The positional reader must be transported with the change of chart, since
the common orbital norm does not make \(\kappa\) invariant under an
isolated renumbering of its components.

\subsection{Depth, dimensions, and physical realizations}
\label{lib06:profundidad}

The bilateral policy that continues one channel and archives the other
has uniform exhaustion of the terminal. Its limit is an
isometry into the sum of all memories, and its inverse is
the convergent adjoint series in
\eqref{lib04:recuperacion-infinita}. A finite number of
complements has two distinct reconstructions: the
truncated adjoint retains an explicit bias, and the least-squares
inverse eliminates it with the computed noise
amplification. The preceding evolutions with a positive terminal
retain that component in the complete archive.

The dimension of the label space, the rank of a projector, and
the dimension of a physical realization are related through
their concrete maps. The unitary on \(1001\) labels and the
twelve-window register are two explicit realizations of
different domains. The exterior conservation law of the preceding
development extends the Gram products to every finite degree,
including the mixed contributions between terminal and memories.
The exceptional and physical realizations remain linked
to their already constructed representations and operators; numerical
ranks alone do not replace that link.

\Needspace{18\baselineskip}
\subsection{Synthesis of conservation under evolution}
The physical actions and the Gaussian protocol of the manuscript
retain the couplings, forms, units, and
preparations involved in their proofs. The bilateral
balance added here extends the collection of conservative
transports and their readouts; its application to a physical
realization is established by composing the operators of that
realization without replacing them with a coincidence of
normalization.

The unity of the development thus resides in a proven
sequence: generation and completion, refinement with memory,
unitary realization, relative symmetry of the reader, and recovery
of structure. The holographic extreme and mean ratio thereby
receives an exposition in which change of scale, incidence,
conservation, and transduction have explicit operations and
jointly retain the information needed to invert
their readouts.
